\begin{lemma}[Change lemma: both vertices outside the starred stage]
Fix a finite loopless multigraph with its exact edge set, a palette of
$k$ colors, and a distinguished edge $e_0$.  Let $(T,\varphi_n)$ be a
positive-rung extended Tashkinov tree.  Let $\mathcal B$ be a hierarchy
base and let
\[
 T_n=T_{n,0}\subset T_{n,1}\subset\cdots\subset
 T_{n,q}\subset T_{n,q+1}=T
\]
be a good hierarchy under $\varphi_n$, where $q$ is its last divider
index and
\[
 A=T_{n,q}^{*}=\noderef{HierarchyStarredStage}(H,q).
\]
Assume that $T$ has maximum property, that $V(T)$ is not elementary,
that $\varphi_n$ is proper and leaves exactly $e_0$ uncolored, and that
the rung number is positive.  These data form a
\noderef{HierarchyCounterexample} $C$.

Assume the six lower-rung assertions of the technical theorem, uniformly
for every positive-rung maximum ETT on the same graph, palette, and
distinguished edge, with arbitrary certified input and generating
colorings and with rung number smaller than that of $T$: lower-rung
elementarity, terminal interchangeability, the two PE stability and
stored-pair assertions, the terminal PE component assertion, the
defective-endpoint order assertion, and the mod-coloring assertion.
Assume also the same-rung induction hypothesis that every maximum ETT
with a good hierarchy having fewer than $q$ dividers is elementary.
Finally assume the large-color estimate
\[
 2n+11\leq |\overline\varphi_n(T_n)|.
\]

Put $\ell=|T|-|A|$, the number of exact tail additions above $A$, and let
\[
 T(y_{\ell-1})=T\big|_{\min\{|A|+(\ell-1),|T|\}}
\]
be the exact ordered prefix supplied by \noderef{TreeSequencePrefix}.
Assume (6.6): $V(T(y_{\ell-1}))$ is elementary under $\varphi_n$.
Assume moreover that $C$ is lexicographically minimal among all
\noderef{HierarchyCounterexample}s on the fixed graph, palette, and
$e_0$ whose ETTs have the same rung number as $T$ and which have exactly
$q$ dividers, allowing their certified input colorings, generating
colorings, bases, ETTs, and good hierarchies to vary.  Explicitly, no
such same-rung counterexample $C'$ satisfies
\[
 \noderef{HierarchyCounterexampleLess}
  \bigl(\noderef{HierarchyCounterexampleMeasure}(C'),
        \noderef{HierarchyCounterexampleMeasure}(C)\bigr).
\]

Let $\alpha\ne\beta$ be colors missing on $T(y_{\ell-1})$.  Write
\[
 i_c=\noderef{HierarchyFirstMissingPosition}(\varphi_n,T,c),
 \qquad v_c=T.\mathit{vertex}(i_c).
\]
If at least one of $\alpha,\beta$ is not missing on $A$, assume both
\[
 v_\alpha\prec v_\beta                                      \tag{a}
\]
and the exact edge-identity exclusion
\[
 \alpha\notin
 \left\{c:\exists t\in\operatorname{Fin}|T|,
       |A|\leq t<i_\beta,
       \ \varphi_n(T.\mathit{edge}(t))=c\right\}.             \tag{b}
\]
Thus (b) is precisely
$\alpha\notin\varphi_n\langle T(v_\beta)-A\rangle$; it is not
an endpoint-only abbreviation.


Assume in addition that both first-missing vertices lie outside the last
starred stage:
\[
        v_\alpha,v_\beta\notin V(A).
\]
Then the full two-color components through $v_\alpha$ and $v_\beta$
are equal in each of the following applicable cases.
\begin{enumerate}
\item $q\geq1$, and either
$\alpha\in\overline\varphi_n(T_{n,q})$ or
$\{\alpha,\beta\}\cap D_{n,q}=\varnothing$.
\item $q=0$, and either
$\alpha\in\overline\varphi_n(T_n)$ or
$\{\alpha,\beta\}\cap D_n=\varnothing$.
\end{enumerate}
\end{lemma}
\begin{proof}
Unpack $C$ and first bind all paper notation.  Put
\[
 H=C.\mathit{hierarchy},\quad
 q=H.\mathit{dividerCount},\quad
 \widehat q=\operatorname{last}(q),\quad
 A=\noderef{HierarchyStarredStage}(H,\widehat q),
\]
put $n=\operatorname{last}(\mathit{ett.rounds})$, and write $T_n$,
$D_n$, $T_{n,q}$, and $D_{n,q}$ for the objects displayed in the
statement.  Put
\[
 i_\alpha=\noderef{HierarchyFirstMissingPosition}(\varphi_n,T,\alpha),
 \quad i_\beta=\noderef{HierarchyFirstMissingPosition}(\varphi_n,T,\beta),
 \quad a=T.\mathit{vertex}(i_\alpha),\quad
 b=T.\mathit{vertex}(i_\beta).
\]
Only now define
\[
 K_x=\noderef{TwoColorComponent}
       (G,\varphi_n,\alpha,\beta,\operatorname{inl}x).
\]
Thus every $K_x$ contains graph-vertex nodes and the individual edge
nodes of the multigraph; no parallel edge identities are collapsed.

The nestedness and terminal fields of
\noderef{TashkinovHierarchy}, together with the two cases in
\noderef{HierarchyStarredStage}, make $A$ an exact ordered initial
segment of $T$.  Let $Y=T(y_{\ell-1})$ denote the prefix in the
statement.  Since $\alpha$ is missing on $Y$, it is missing at some
displayed vertex of $Y$.  The least-candidate definition of $i_\alpha$
then puts $a$ in $V(Y)$ and makes $\alpha$ missing at $a$.  The same
argument puts $b$ in $V(Y)$ and makes $\beta$ missing at $b$.  Here we
used \noderef{TreeSequencePrefix} and the entrywise relation
\noderef{TreeSequenceInitialSegment}: every missing-color candidate in
$Y$ is the same candidate at the same index in $T$.

The location hypothesis says $a,b\notin V(A)$.  Since $A$ and $Y$ are
both exact prefixes of $T$ and $b\in V(Y)-V(A)$, their entrywise orders
give $A\prec Y$.  In particular $\ell\geq2$: if $\ell\leq1$, then
$\ell-1=0$ and the defining length
$\min\{|A|+(\ell-1),|T|\}$ of $Y$ is at most $|A|$.  This contradicts
$A\prec Y$.  Every starred and ordinary hierarchy stage through
$T_{n,q}$ is an initial segment of $A$, hence of $Y$.  Restricting the
pairwise-disjoint condition of
\noderef{ColoredVertexSetElementary} from $V(Y)$ therefore makes each of
those stages elementary.

Neither $\alpha$ nor $\beta$ is missing on $A$.  For example, if
$\alpha$ were missing at $x\in V(A)$, then $x\ne a$ because
$a\notin V(A)$, while both $x$ and $a$ lie in the elementary set $V(Y)$
and both miss $\alpha$, a contradiction.  The argument for $\beta$ is
identical.  Thus the premise of (a)--(b) holds.  If $a=b$, then
$K_a=K_b$ definitionally and both required implications are finished.
Henceforth suppose $a\ne b$.  Applying (a)--(b), we obtain
\[
                         a\prec_T b                              \tag{1}
\]
and the stated exact exclusion of $\alpha$ from all displayed edges
between $A$ and the prefix through $b$.  Equation (1) is the order used
below.

Define the exact ordered prefix
\[
 B=T(b)=\noderef{TreeSequencePrefix}(T,i_\beta).
\]
It ends at the actual vertex entry $b$.  Equation (1), the two location
hypotheses, and entrywise prefix comparison give
\[
 T_n\preccurlyeq A\prec B\preccurlyeq Y\prec T.                \tag{2}
\]
The last strict relation follows also by lengths: if
$s=|B|-|A|$, then $b\in V(Y)$ gives
\[
                  1\leq s\leq\ell-1<\ell.                     \tag{3}
\]

Before changing the coloring, we construct all old-color data on this
exact prefix.  Keep the entire certified series of the given ETT and
replace only its displayed terminal tree $T$ by $B$.  The last ladder
tree $T_n$ is a strict ordered prefix of $B$ by (2), and $B$ is an
ordered prefix of $T$, hence of the final certified closure.  These are
exactly the two terminal-tree fields of
\noderef{ExtendedTashkinovTree}; they give an old-color ETT
$\mathcal E_B$ on $B$ with the same positive rung number and the same
generating coloring $\varphi_n$.  Since
\noderef{MaximumProperty} depends only on that rung and certified series,
the maximum-property field of $C$ is also a maximum-property witness for
$\mathcal E_B$.

Copy the core, closure, and support position of the old
\noderef{HierarchyBase}.  All its conditions refer to the unchanged
terminal ladder tree, terminal kind, generating series, and those copied
objects, so they form a base $\mathcal B_B$ for $\mathcal E_B$.  Define a
hierarchy $H_B$ with the same divider count $q$, the same stages through
$T_{n,q}$, the same divider positions, and terminal stage $B$.  The only
new nesting and strictness assertion is $T_{n,q}\prec B$; it follows from
$T_{n,q}\preccurlyeq A\prec B$.  Every divider-stage equality remains an
equality of exact prefixes because both its old stage and $B$ are
entrywise prefixes of $T$.  The core-inside-first-stage condition is
copied when $q>0$, while for $q=0$ it is exactly
$\mathcal B_B.\mathit{core}=A\prec B$.  This verifies every field of
\noderef{TashkinovHierarchy} for $H_B$.

We also construct, still before recoloring, a genuine good hierarchy on
$H_B$.  Use the old Gamma function.  The starred stages, ordinary stages,
certified stored set, and their missing-color sets are unchanged, so the
sets in \noderef{HierarchyStoredColorsAt} and
\noderef{HierarchyStageMissingColors} are unchanged at every index.  For
the avoidance field, fix $j,\eta$ and an exact edge position in the
$(j,\eta)$ segment of $H_B$.  If an index is smaller than the first
$\eta$-missing position in $B$, it is also smaller than the first
$\eta$-missing position in $T$: otherwise an earlier candidate in $T$
would be the same vertex entry of the prefix $B$, contradicting
minimality there.  Consequently the new segment is contained, edge by
edge and with the same $\varphi_n$-values, in the old segment defined by
\noderef{HierarchySegmentColors}.  Old avoidance therefore implies the
required avoidance for $H_B$.  Pair size and pairwise disjointness copy
literally; unique refresh copies because its two stages and their vertex
difference are unchanged; the PE base field copies because the series,
base, zeroth stored set, and zeroth Gamma family are unchanged; and every
minus-closed field copies because it concerns an unchanged ordinary
stage.  These are all five fields of \noderef{GoodHierarchy}; denote the
resulting witness by $\mathcal G_B$.

We now prove either implication in the statement.  Assume its antecedent.
If $q>0$, then $A=T_{n,q}$ by
\noderef{HierarchyStarredStage}.  The first disjunct of alternative 1 is
false because $\alpha$ is not missing on $A$, so its second disjunct gives
$\{\alpha,\beta\}\cap D_{n,q}=\varnothing$.  If $q=0$, then $T_n$ is an
initial segment of $A$, so neither color is missing on $T_n$; the first
disjunct of alternative 2 is false and its second disjunct gives
$\{\alpha,\beta\}\cap D_n=\varnothing$.  In the positive-$q$ case,
\noderef{HierarchyStoredColorsAt} says
\[
 D_{n,q}=D_n\setminus\overline\varphi_n(A).
\]
Since neither color is in the subtracted set, exclusion from $D_{n,q}$
also excludes both from $D_n$.  Thus in both alternatives
\[
 \{\alpha,\beta\}\cap
   \bigl(D_n\cup\overline\varphi_n(A)\bigr)=\varnothing.       \tag{4}
\]

We next exclude both colors from every displayed edge of $B$.  The
terminal constructor in \noderef{TashkinovSeriesStep} has one of three
kinds.  In the PE case its final closure supplies a
$\varphi_n$-\noderef{TashkinovAugmentationSequence} from $T_n$ to the
final certified closure.  In the SE or RE case it supplies an exact edge
$f_n$, the equality $F_n=\{f_n\}$, the exact seed $T_n+f_n$, and such an
augmentation sequence from that seed.  In either case the seed is an
initial segment of $B$: for PE this follows from $T_n\prec B$; for SE/RE,
the seed and $B$ are prefixes of the same final closure and
$|B|>|T_n|$, so the one-edge seed has length at most $|B|$.  Apply
\noderef{TashkinovAugmentationSequenceRestrict} to the seed, $B$, and the
final closure.  We obtain the corresponding exact old-color augmentation
sequence ending at $B$.                                             \tag{5}

By \noderef{HkuTreeColorsControlled}, every nonroot displayed color on
$T_n$ lies in
$\overline\varphi_n(T_n)\cup D_n$, a subset of
$\overline\varphi_n(A)\cup D_n$; hence (4) excludes both colors there.
The root identity is $e_0$ and is uncolored by
\noderef{IsOneEdgeUncolored}.  In SE, the selected
\noderef{MaximumDefectiveChoice} gives
$\varphi_n(f_n)=\delta_n$ and the constructor puts
$\delta_n\in S_n$.  In RE, the
\noderef{RevisitingExtensionCertificate} puts $f_n$ on the displayed
$(\gamma_h,\delta_h)$-cycle, and the constructor puts both cycle colors
in $S_n$.  Thus in either non-PE case the color of $f_n$ belongs to the
raw stored union.  By \noderef{TashkinovStoredColors}, it therefore lies
in $\overline\varphi_n(T_n)\cup D_n$, so (4) excludes both colors from
$f_n$ as well.

Every remaining edge in (5) is a one-step
\noderef{TashkinovAugmentation}: its exact new edge color is missing at a
vertex of its old prefix.  If the new edge occurs no later than $A$, that
witness is a vertex of $A$, so (4) excludes both $\alpha$ and $\beta$.
If it occurs strictly after $A$ and no later than the edge that introduces
$b$, its old-prefix witness is strictly before $b$.  It cannot miss
$\beta$, by the least-candidate definition of $i_\beta$, and the exact
edge-identity exclusion (b) forbids its new edge from having color
$\alpha$.  These cases exhaust (5), including its last edge.  Therefore
\[
  \varphi_n(g)\notin\{\alpha,\beta\}
  \quad\hbox{for every colored displayed edge identity }g\hbox{ of }B.
                                                                    \tag{6}
\]

Suppose for contradiction that $K_a\ne K_b$, and set
\[
             \sigma_n=\varphi_n/K_b.
\]
By \noderef{TwoColorComponentsDisjoint}, the two components are disjoint,
so the component through $a$ is untouched.  The definition
\noderef{TwoColorComponentSwap} changes only edge nodes of $K_b$ and only
transposes $\alpha$ and $\beta$.  It fixes every uncolored value;
\noderef{TwoColorComponentSwapProper} makes $\sigma_n$ proper, and
\noderef{IsOneEdgeUncolored} then shows that exactly $e_0$ remains
uncolored.  The color $\alpha$ is still missing at $a$.  At $b$, the old
color $\beta$ is missing, while elementarity of $Y$, $a\ne b$, and the
fact that $\alpha$ is missing at $a$ show that $\alpha$ is present.
This is the complement relation in \noderef{MissingColors} and
\noderef{PresentColors}.  Properness gives its unique incident
$\alpha$-edge; that edge node lies in
$K_b$ and is changed to $\beta$.  Hence $\alpha$ is missing at $b$ under
$\sigma_n$.                                                     \tag{7}

We need a pointwise statement stronger than stability.  Let $x$ be any
vertex of $B$ strictly before $b$.  If its vertex node is outside $K_b$,
then no incident edge node can lie in $K_b$, and every incident value is
fixed.  If its vertex node lies in $K_b$, then $x\ne a$, since otherwise
$K_a=K_b$.  The first-missing property of $b$ says that $\beta$ is
present at $x$.  Elementarity of $Y$, applied to the distinct vertices
$x,a$, says that $\alpha$ is also present at $x$.  Properness gives
exactly one incident edge of each color, both edge nodes lie in $K_b$,
and the swap exchanges their values.  All other incident values are
fixed.  Thus in both cases
\[
       \overline{\sigma_n}(x)=\overline\varphi_n(x)
       \qquad(x\in V(B)-\{b\}).                                \tag{8}
\]
At $b$, the same inspection shows that missing membership is unchanged
for every color outside $\{\alpha,\beta\}$.  Equation (6) and the
definition of the swap also give, for every displayed edge identity of
$B$,
\[
                         \sigma_n(g)=\varphi_n(g).              \tag{9}
\]

Equation (8) holds on $T_n$ and on every hierarchy stage.  By (4), neither
changed color belongs to $D_n$ or is missing on $T_n$.  Hence every edge
whose old color is controlled by
\noderef{StableColoring} is fixed, and (8) supplies its vertex clause.
Therefore $\sigma_n$ is $(T_n,D_n,\varphi_n)$-stable.  In the PE case,
the ordered-union clause of the base puts every vertex of $R_n$ in the
core and hence in $A$.  Neither changed color is missing on $R_n$, so the
same argument proves $(R_n,\varnothing,\varphi_n)$-stability.  Together
these two relations are exactly the
\noderef{CombinedStableColoring} condition.                      \tag{10}

Transport the restricted sequence (5) one augmentation at a time.  Its
new exact edge identity is displayed in $B$, so (9) preserves its color.
Its missing-color witness lies in the preceding prefix, hence strictly
before the final vertex $b$, so (8) preserves that witness.  The length,
attachment, and exact initial-segment clauses are independent of the
coloring.  Thus every old step is a $\sigma_n$-augmentation.  We obtain
a $\sigma_n$-augmentation sequence from $T_n$ to $B$ in PE, and from the
same exact seed $T_n+f_n$ to $B$ in SE or RE.                    \tag{11}

We now transport the certified ETT, keeping the extension kinds separate.
Suppose first that the terminal kind is PE.  Apply
\noderef{CombinedStableClosurePreserved} to $\mathcal E_B$,
$\mathcal B_B$, and (10).  It proves that the same ordered $R_n$ is still
a closure of the same supporting prefix under $\sigma_n$ and that all
displayed identities of $T_n$ and $R_n$ retain their values.  Apply
\noderef{HierarchySplitter} to the old ETT $\mathcal E_B$, its maximum
property, the given lower-rung package, the old hierarchy $H_B$, (10),
and the PE sequence in (11).  Its conclusion supplies an actual
\noderef{CorrespondingExtendedTashkinovTree} $\mathcal E'_B$ on the exact
ordered tree $B$, maximum property for that replacement ETT, a base
$\mathcal B'_B$, a hierarchy $H'_B$, and
\[
             \noderef{SameHierarchyStages}(H_B,H'_B).          \tag{12}
\]

Suppose instead that the terminal kind is SE or RE.  The first relation
in (10) and \noderef{Td}, applied to $\mathcal E_B$, its maximum
property, positivity, and the lower-rung package, prove that $\sigma_n$
is a \noderef{ModColoring} of $\mathcal E_B$.  The exact connecting edge,
singleton equality $F_n=\{f_n\}$, exact one-edge seed, and the non-PE
sequence in (11) make $B$ eligible for the universal clause of that mod
coloring.  It supplies an actual
\noderef{CorrespondingExtendedTashkinovTree} $\mathcal E'_B$ on $B$.
Then \noderef{MaximumPropertyTransport} supplies maximum property for
$\mathcal E'_B$.

Build $\mathcal B'_B$ with core $T_n$ and copy the unused closure and
support position.  Its displayed coloring and the replacement ETT's
generating coloring are both $\sigma_n$, so the terminal missing-color
compatibility is reflexive; the parallel implication is vacuous, and the
nonparallel field is the equality of the core with $T_n$.  Copy the
old stages through $T_{n,q}$, the divider positions, and terminal stage
$B$ to obtain $H'_B$.  The verification is the same entrywise verification
used to construct $H_B$.  The correspondence certificate identifies the
terminal ladder tree and its non-PE kind, so the copied data satisfy every
clause of \noderef{SameHierarchyStages}; thus (12) also holds in this
case.

It remains to construct goodness under $\sigma_n$; it is not a consequence
of (12).  Transport the Gamma function of $\mathcal G_B$ across the divider
count equality in (12).  The corresponding-ETT certificate identifies
every raw stored set.  Every starred stage lies in $A$, so (8) identifies
its missing-color set vertex by vertex.  Unfolding
\noderef{HierarchyStoredColorsAt} therefore gives, at every stage,
\[
                         D'_{n,j}=D_{n,j}.                      \tag{13}
\]
The same argument on $T_n$ and the ordinary stages gives equality of all
sets in \noderef{HierarchyStageMissingColors}.  Pair cardinality and
pairwise disjointness for the transported Gamma pairs now copy literally
from $\mathcal G_B$.

For avoidance, fix $j$ and $\eta\in D'_{n,j}=D_{n,j}$.  The definition
\noderef{HierarchyStoredColorsAt} gives $D_{n,j}\subseteq D_n$; by (4),
$\eta$ is neither $\alpha$ nor $\beta$.  Equation (8), and the final
observation after it at $b$, show pointwise that $\eta$ is missing at a
vertex of $B$ under $\sigma_n$ exactly when it is missing there under
$\varphi_n$.  Hence the two first $\eta$-missing positions in the exact
tree $B$ agree.  Now take an individual edge position in the new
$(j,\eta)$ segment.  Equation (12) identifies its starred and next stages
and (9) identifies the color of that exact edge.  The equal first-missing
position makes it the corresponding edge of the old $(j,\eta)$ segment,
where the old Gamma pair avoids its color.  This proves avoidance
pointwise at every paper-relevant $(j,\eta)$; no equality of the full
\noderef{HierarchySegmentColors} functions is asserted.

The unique-refresh field also copies: (12) fixes both stage differences,
(13) fixes the stored sets, (8) fixes every missing-color set on those
differences, and the Gamma pairs are the transported old pairs.  We next
check the PE base field.  In that case
\noderef{CombinedStableClosurePreserved} kept the exact same closure,
while (12) kept the exact same core; (8) identifies missing sets at every
vertex of both.  Hence
the two common-missing sets, the zeroth stored set, and the zeroth Gamma
union are exactly the old ones, so the required cardinality is unchanged.
For non-minus-closedness, unpack its old witness: there are a color $c$
missing on the core but outside the zeroth Gamma union and an exact
boundary edge $g$ with $\varphi_n(g)=c$.  Neither $\alpha$ nor $\beta$
is missing on the core, because its vertices lie in $A$; hence
$c\notin\{\alpha,\beta\}$.  The swap fixes $g$, (8) preserves the missing
witness, and the boundary relation is structural.  The same $c,g$ witness
non-minus-closedness under $\sigma_n$.

Finally fix an internal stage and a color $c$ to which its old
minus-closed field must be transferred.  Thus $c$ is missing on that
stage under $\sigma_n$ and is outside the preceding Gamma union.
Equation (8) makes it missing there under $\varphi_n$ as well.  The stage
is contained in $A$, so (4) implies
$c\notin\{\alpha,\beta\}$.  Any exact boundary edge of color $c$ under
$\sigma_n$ is therefore fixed by the swap and has color $c$ under
$\varphi_n$, contradicting the old field.  This proves every
minus-closed clause in the definition
\noderef{TreeColorSetMinusClosed}.  We have now checked all five fields of
\noderef{GoodHierarchy}; call the resulting witness $\mathcal G'_B$.
Only at this point do (12), (13), and equality of the transported Gamma
pairs give
\[
       \noderef{SameGoodHierarchy}(\mathcal G_B,\mathcal G'_B). \tag{14}
\]

By (7), the distinct vertices $a,b\in V(B)$ both miss $\alpha$ under
$\sigma_n$, so $V(B)$ is not elementary.  Positivity is unchanged by the
corresponding ETT.  Together with the maximum property, base, hierarchy,
and actual good witness just constructed, these facts form a
\noderef{HierarchyCounterexample} $C'$ on $B$, and (12) gives it exactly
$q$ dividers.  The old prefix ETT $\mathcal E_B$ was defined with a
literal copy of the certified series of the original ETT, so their rung
numbers are equal.  The `rungAgreement' field of the
\noderef{CorrespondingExtendedTashkinovTree} certificate then gives
$(\mathcal E'_B).\mathit{rounds}=\mathcal E_B.\mathit{rounds}$.
Composing these two exported equalities proves that the ETT underlying
$C'$ has the same rung number as the ETT underlying $C$.

It remains to compare measures.  Put
\[
 r=\noderef{HierarchyPathNumber}(H),\qquad
 r'=\noderef{HierarchyPathNumber}(H'_B),
\]
and retain $s$ from (3).  If $s<r$, the fallback-and-candidate definition
of the path number gives $r'\leq s<r$.  If $r\leq s$, then the original
tail beginning at its $r$-th added vertex is a path.  Indeed, the fallback
value $r=\ell$ is impossible because $r\leq s<\ell$; hence $r$ is an
actual candidate in the definition.  Its exact initial segment ending in
$B$ is again a path, so $r$ is a candidate for $H'_B$ and $r'\leq r$.
Thus either $r'<r$, or $r'=r$.  In the latter case the ordinary last
stage $T_{n,q}$ is unchanged while $B\prec T$, so
\[
 |B|-|T_{n,q}|<|T|-|T_{n,q}|.
\]
By \noderef{HierarchyCounterexampleMeasure} and
\noderef{HierarchyCounterexampleLess}, in both cases
\[
 \noderef{HierarchyCounterexampleMeasure}(C')
 <_{\mathrm{hier}}
 \noderef{HierarchyCounterexampleMeasure}(C).                  \tag{15}
\]
The just-proved rung equality and the divider-count equality make $C'$
eligible for the stated comparison class, so (15) contradicts the
assumed minimality.
Therefore $K_a=K_b$ for the chosen applicable alternative.  Applying the
same argument separately to alternatives 1 and 2 proves both implications
in the statement.
\end{proof}
